\section{The algorithm}
\label{sec:algorithm}

\subsection{Field and ring choices}
\label{sec:field-choices}

Fix
\[
  f(X) \;=\; X^{16} + X^{5} + X^{3} + X^{2} + 1 \;\in\; \F_2[X],
\]
an irreducible polynomial of degree~$16$. Define
$\GF(2^{16}) = \F_2[X]/f$ and $\flift \in \Z[X]$ to be the integer
polynomial with the same set of monomials and the same $\{0,1\}$
coefficients. The lifted ring is
\[
  \Rlift \;=\; \Z[X]/\flift,
\]
a $16$-dimensional $\Z$-module on which addition is component-wise and
multiplication is schoolbook poly-multiplication followed by reduction
modulo $\flift$. The reduction is
\[
  X^{16} \;\equiv\; X^{5} + X^{3} + X^{2} + 1 \pmod{\flift},
\]
so any product whose temporary degree exceeds $15$ is folded back by
contributing to four low positions per high monomial.

\subsection{The Cantor basis}
\label{sec:cantor}

The Cantor basis of $\GF(2^{16})$ is the sequence
$\bb_0, \bb_1, \ldots, \bb_{15}$ defined recursively by
\begin{equation}
  \bb_0 = 1, \qquad \bb_i^2 + \bb_i = \bb_{i-1} \quad \text{for } i \ge 1.
  \label{eq:cantor-recurrence}
\end{equation}
The recurrence is solvable because $\Tr(\bb_{i-1}) = 0$ at every step
($\Tr(\bb_i^2 + \bb_i) = \Tr(\bb_i)^2 + \Tr(\bb_i) = 0$ in $\F_2$).
For~$i = 0$, $\Tr(1) = 16 \cdot 1 = 0$ since the extension degree is
even.

The two crucial properties for the additive FFT are:
\begin{itemize}
  \item $V_i := \mathrm{span}_{\F_2}(\bb_0, \ldots, \bb_{i-1})$ is a
    chain of $\F_2$-subspaces of $\GF(2^{16})$, with
    $\dim V_i = i$.
  \item The associated subspace polynomial
    \[
      s_i(X) \;=\; \prod_{\alpha \in V_i} (X + \alpha)
    \]
    is $\F_2$-linearised: only monomials of degree $2^j$ for
    $j = 0, \ldots, i$ appear, all coefficients lie in $\F_2$, and
    the leading coefficient is~$1$. The Cantor basis additionally
    guarantees $s_i(\bb_i) = 1$, which simplifies the
    recurrence to
    \begin{equation}
      s_{i+1}(X) \;=\; s_i(X)^2 + s_i(X).
      \label{eq:cantor-subspace-recurrence}
    \end{equation}
\end{itemize}

Crucially for the structural lift: every $s_i$ coefficient is in
$\F_2$, so its lift to $\Z$ is in $\{0, 1\}$. Combined with the fact
that Cantor basis vectors themselves are elements of $\GF(2^{16})$,
every twiddle the FFT uses --- and hence every entry of the radix-8
Vandermonde matrices below --- is a $\{0, 1\}$-coefficient polynomial
once lifted to $\Rlift$.

\subsection{The novel polynomial basis}
\label{sec:lch}

Let $m = \log_2 N$, where $N$ is a power of $2$ at most $2^{16}$.
Express the input polynomial $P(X)$ of degree~$< N$ in the
\emph{novel polynomial basis}
\[
  X_i(X) \;=\; \prod_{j \in \{0..m-1\}\,:\, i_j = 1} s_j(X),
\]
where $i_j$ is the $j$-th bit of $i$. Then $P = \sum_{i=0}^{N-1} c_i \, X_i$,
and the forward FFT maps the coefficient vector
$(c_0, \ldots, c_{N-1})$ to the evaluation vector $(P(v_0), \ldots,
P(v_{N-1}))$ at the points $v_k = \sum_l k_l \bb_l \in V_m$. Data is
in natural order at both input and output --- no bit-reversal
permutation is required.

The textbook radix-2 LCH FFT computes this in $m$ stages: at stage $i$,
each block of $2^{i+1}$ entries is processed by $2^i$ butterflies, each
mixing a pair $(a,b)$ via a single twiddle $t = s_i(\omega)$:
$a' = a + t\cdot b$, $b' = a' + b$. We do \emph{not} use the radix-2
form directly --- see below.

\subsection{The radix-8 evaluator}
\label{sec:radix8}

The shipped encoder uses a \textbf{radix-8} variant. Three consecutive
radix-2 stages are merged into one \emph{meta-stage}: instead of
$m$ stages of $2\times 2$ butterflies, the FFT runs $m/3$ meta-stages,
each applying an $8 \times 8$ Vandermonde matrix per block of $8$
codeword positions. (This requires $3 \mid m$; the linear-code
constructor rejects parameters that violate it.)

Concretely, meta-stage $K$ covers the three subspace polynomials
$s_{3K}, s_{3K+1}, s_{3K+2}$. For each block $B$ of $8$ positions
the $8 \times 8$ matrix $V_{K,B}$ has entries
\[
  V_{K,B}[j][i] \;=\; X^{(K)}_i\!\left( v_{K,B,j} \right)
  \;=\;
  \prod_{a \,:\, i_a = 1} s_{3K+a}\!\left( v_{K,B,j} \right),
\]
the $i$-th novel-basis polynomial of the $3$-dimensional
sub-recursion, evaluated at the $j$-th point of that block. The
meta-stage output is the per-block matrix--vector product
\[
  \mathrm{out}[j] \;=\; \sum_{i=0}^{7} V_{K,B}[j][i] \cdot \mathrm{in}[i].
\]
Composed over the $m/3$ meta-stages, this evaluates exactly the same
multilinear extension as the radix-2 FFT --- over $\GF(2^{16})$ the two
agree identically.

\paragraph{Why radix-8, and why it matters for the lift.}
The radix-8 form is the load-bearing design choice for this code,
\emph{not} a peephole optimization. The reason is the magnitude of the
lifted codeword coefficients (\autoref{sec:optimization}):

\begin{itemize}
  \item In a radix-2 $\Z$-lift, the codeword is computed as a chain of
    $m$ sequential multiplications in $\Rlift$ by lifted twiddles. Each
    multiplication is a $\Z$-arithmetic \code{mul + reduce-mod-}$\flift$,
    so $\Z$-carries accumulate across all $m$ stages. Empirically the
    coefficients reach $\sim 2^{47}$ at $\code{codeword\_len} = 4096$.
  \item In the radix-8 form, each $V_{K,B}$ is computed \emph{in
    $\GF(2^{16})$} --- where the product of three subspace-polynomial
    values is just another $\GF(2^{16})$ element --- and only then
    lifted, so every Vandermonde entry has $\{0,1\}$ coefficients. The
    FFT then composes $m/3$ such $\{0,1\}$-coefficient matvecs in
    $\Rlift$. With only $m/3$ carry-accumulating meta-stages instead of
    $m$, and a $\{0,1\}$-coefficient multiplier at each, the codeword
    coefficients stay near $2^{26}$ ($\sim 28$ signed bits) --- a
    $\sim 21$-bit reduction over the radix-2 $\Z$-lift.
\end{itemize}

This is a genuine change of code: the radix-8-with-GF-lifted-Vandermonde
encoder produces a \emph{different} lifted code than a radix-2 $\Z$-lift
would. The two agree modulo $2$ (both reduce to the standard
$\GF(2^{16})$ RS codeword, see below), but differ in the higher
$\Z$-bits, because precomputing $V_{K,B}$ in $\GF(2^{16})$ discards the
inter-stage $\Z$-carries that a sequential radix-2 lift would keep.
Choosing the radix-8 code is choosing the variant with the small
coefficients.

\subsection{The structural lift and mod-2 commutativity}
\label{sec:lift}

The lifted FFT runs the radix-8 algorithm with arithmetic in $\Rlift$
rather than $\GF(2^{16})$. The Cantor basis, subspace polynomials, and
the $8 \times 8$ Vandermonde matrices are all computed once in
$\GF(2^{16})$ and then \emph{lifted} coefficient-by-coefficient from
$\F_2$ to $\{0, 1\} \subset \Z$. The inputs are \code{BinaryPoly<16>}
elements, whose $\F_2$-coefficients also lift trivially.

Each per-block matvec is a sequence of $\Rlift$-additions and
multiplications by lifted, $\{0,1\}$-coefficient Vandermonde entries.
The multiplication is a \emph{sparse shift-and-accumulate}: for every
set bit of the entry, add the operand shifted by that many positions
to a length-$31$ unreduced accumulator, then fold the accumulator back
to length~$16$ via the $\flift$ reduction.

The soundness witness is the commutative diagram: let
$\pi_2 : \Rlift \to \GF(2^{16})$ be coefficient-wise mod-$2$ reduction.
Then for any input vector $\mathbf c$,
\[
  \pi_2\!\bigl( \mathrm{LiftedRadix8FFT}(\mathrm{lift}(\mathbf c)) \bigr)
  \;=\;
  \mathrm{FFT}_{\GF(2^{16})}(\mathbf c).
\]
This holds because mod-$2$ reduction is a ring homomorphism that
commutes with addition and the $\flift$ reduction (as
$f \equiv \flift \pmod 2$), the lifted Vandermonde entries reduce mod
$2$ to the $\GF(2^{16})$ entries they were lifted from, and the
$\GF(2^{16})$ radix-8 evaluator computes exactly the standard
$\GF(2^{16})$ additive-FFT codeword. So the lifted code reduces mod $2$
to a genuine $\GF(2^{16})$ Reed--Solomon code, which is the basis for
inheriting that code's distance.
